\documentclass[12pt,a4paper,oneside]{article}
\usepackage[brazil]{babel}
\usepackage[utf8]{inputenc}
\usepackage{olymp}

\setcounter{problem}{0}

\begin{document}

\begin{problem}{Amplitude térmica}{amplitude}{2}

Os últimos dias têm sido bem frios em Viçosa. Em algumas dias a temperatura aumenta após as primeiras horas da manhã, mas em outros dias ela permanece baixa o dia todo.

Nesse exercício estamos interessados em estudar a amplitude térmica na cidade. A amplitude térmica de um certo dia é a diferença entre a maior e a menor temperatura registrada naquele dia.
Como exemplo, veja a tabela abaixo, que mostra os dados de Viçosa na semana passada. No dia 17, a máxima temperatura foi de $27 ^\circ$C, enquanto a mínima foi de $17 ^\circ$C. Então, a amplitude térmica nesse dia foi de $27 - 17 = 10 ^\circ$C.


\begin{table}[h]
\centering
\begin{tabular}{|c|c|c|c|c|c|c|c|}
\hline
Dia (maio/2026)      & 17 & 18 & 19 & 20 & 21 & 22 & 23 \\
\hline
Temperatura Máxima ($^\circ$C) & 27 & 26 & 23 & 21 & 20 & 19 & 24 \\
Temperatura Mínima ($^\circ$C) & 17 & 17 & 17 & 16 & 14 & 13 & 14 \\
Amplitude térmica ($^\circ$C) & 10 & 9 & 6 & 5 & 6 & 6 & 10 \\
\hline
\end{tabular}
%\caption{Daily maximum and minimum temperatures with thermal amplitude}
\label{tab:temps}
\end{table}

Faça um programa para calcular a amplitude térmica dos dias informados na entrada.

%\Notes
%
%Observações, se tiver.

\InputFile

A entrada contém três linhas: a primeira contém um número inteiro $N$, indicando a quantidade de dias em estudo;
a segunda contém as temperaturas máximas em cada um dos $N$ dias;
e a terceira as temperaturas mínimas em cada um dos $N$ dias.

Restrições: $1 \leq N \leq 1000$, $-60 \leq \text{temperaturas} \leq 60$. Todos os valores da entrada são números inteiros.

\OutputFile

Seu programa deve escrever uma linha com a amplitude térmica de cada dia.

%\Notes

\Example

\begin{example}
\exmp{
7
27 26 23 21 20 19 24
17 17 17 16 14 13 14
}{
10 9 6 5 6 6 10
}%
\end{example}

\begin{example}
\exmp{
3
20 20 20
15 20 18
}{
5 0 2
}%
\end{example}

%Comentários sobre o exemplo, se necessário.

\end{problem}

\end{document}